\subsection{Testing the Independence Explanation}
\label{sec:independence}

If correlated errors are the obstacle, then making the reviewers more independent should reduce false consensus even without changing the gate.
Two things could make the reviewers of Section~\ref{sec:pilot} non-independent: they share the author's backbone, and their prompt contains the author's draft (with the instruction to solve the problem first).
We varied both, on the same development split and with the same agreement gate: \emph{blind} reviewers never see the draft and simply solve the problem; \emph{heterogeneous} reviewers use a different backbone (GPT-4o-mini) while the author keeps GPT-4.1-mini; and a variant combines both.
Each variant was run once, and the blind same-backbone variant a second time (Table~\ref{tab:independence}).

\begin{table}[!t]
\caption{Independence interventions on the development split (210 questions, author GPT-4.1-mini; one run each, two for the blind variant). Reproduce: share of incorrect drafts for which \emph{both} reviewers independently give the author's wrong answer. Rev.\ acc.: accuracy of the reviewers' own first-round answers. Other columns as in Table~\ref{tab:gate}.}
\label{tab:independence}
\footnotesize
\setlength{\tabcolsep}{2.5pt}
\resizebox{\linewidth}{!}{\begin{tabular}{l cc c cc cc c c r}
\toprule
 & \multicolumn{2}{c}{Reviewers} & & \multicolumn{2}{c}{Pass in round 1} & & & & & \\
\cmidrule(lr){2-3}\cmidrule(lr){5-6}
Variant & see draft & backbone & Rev.\ acc. & Correct & Incorr. & Reproduce & Repaired & Net & Lost & Tokens \\
\midrule
Agreement gate (Sec.~\ref{sec:pilot}) & yes & same & 84 & 98 & 80 & 80 (28/35) & 71 (5/7) & $+1.9$ & 0.5 & 2.3k \\
\quad blind reviewers, run 1 & no & same & 83 & 95 & 48 & 48 (19/40) & 43 (9/21) & $+2.4$ & 1.9 & 2.0k \\
\quad blind reviewers, run 2 & no & same & -- & 95 & 61 & 61 (23/38) & 40 (6/15) & $+1.4$ & 1.4 & 1.7k \\
\quad heterogeneous reviewers & yes & different & 81 & 97 & 84 & 84 (31/37) & 17 (1/6) & $-0.5$ & 1.0 & 2.4k \\
\quad blind + heterogeneous & no & different & 78 & 84 & 42 & 42 (14/33) & 32 (6/19) & $+0.5$ & 2.4 & 2.3k \\
\bottomrule
\end{tabular}}
\end{table}

\paragraph{Exposure to the draft, not the shared backbone, is the dominant measured source of correlation.}
Hiding the draft from same-backbone reviewers reduces the share of incorrect drafts that both reviewers reproduce from 80\% to 48\% in one run and 61\% in a second, and reduces the share that pass the gate by the same amounts, while the reviewers' own accuracy is unchanged (83\% versus 84\%).
Replacing the backbone of the reviewers while still showing them the draft does nothing: a different and somewhat weaker model that reads the author's answer reproduces it in 84\% of the author's failures, if anything more often than the author's own backbone does.
Combining both interventions gives the lowest reproduction rate (42\%) but no clear additional benefit over blind same-backbone reviewers (48--61\%), and the weaker reviewers challenge more correct drafts (16\% versus 5\%).
Reviewer skill is not the explanation: the four reviewer configurations solve the problems with similar accuracy (78--84\%) when their answers are taken on their own.
Our results therefore indicate that, in this setting, exposure to the candidate answer rather than model identity is the dominant source of correlated reviewer errors; we do not claim that this holds for every multi-agent system, but it does show that model heterogeneity alone does not amount to independent judgment.

\paragraph{Where the bottleneck sits: entry into reconsideration, not repair.}
Reading Tables~\ref{tab:pilot} and \ref{tab:independence} together as the three stages of Section~\ref{sec:agreement} locates the failure precisely.
Once an incorrect draft \emph{does} enter reconsideration under the agreement gate, repair is comparatively effective: 5 of 7 such drafts are fixed with draft-exposed reviewers and 9 of 21 with blind reviewers, against 5 of 27 under the risk-based gate, whose challenges carry no alternative solution.
The dominant loss is upstream: with draft-exposed reviewers 80\% of incorrect drafts never enter reconsideration at all, because the reviewers confirm them.
Blind review moves the collective from the first bottleneck to the second: it roughly doubles the share of errors that are challenged (from 20\% to 39--52\% across the two runs), at a cost \emph{below} that of the ungated pipeline (1.7--2.0k versus 3.4k tokens), and the remaining limit is the repair stage, where 40--43\% of challenged errors are fixed in both runs.
The within-run gains (\(+2.4\) and \(+1.4\) points) are within the run-to-run variation documented in Table~\ref{tab:repcons}, and we draw no accuracy conclusion from 210 questions; what the experiment establishes, in both runs, is the mechanism.

\paragraph{What remains.}
The residual 42--61\% overlap of errors under blind review reflects failures that arise even without exposure to the author's draft: for these questions, independent solvers of the same or a similar model make the same mistake.
Such errors are unlikely to be resolved by blinding alone, and they motivate stronger sources of solver diversity---different model families, tools or retrieval, or different reasoning representations---together with an external check that does not rely on agreement at all.
